\clearpage
\VerdierPhysicalPageAnchor{0248}%
\begingroup\small\setlength{\parskip}{0pt}\setlength{\emergencystretch}{1.5em}
\setlength{\abovedisplayskip}{4pt plus 2pt minus 2pt}%
\setlength{\belowdisplayskip}{4pt plus 2pt minus 2pt}%

\VerdierSemanticAnchor{ch4d-4-6-8}
\noindent\textbf{Proposition 4.6.8.}\quad
\textit{Let $Y$ be a complex of $\mathcal A$, and let
$\varepsilon\colon Y\longrightarrow I^{\bullet\bullet}$ be a
Cartan--Eilenberg injective resolution of $Y$. When $Y$ is an object of
$\operatorname{comp}^{+}(\mathcal A)$, or when the resolution
$I^{\bullet\bullet}$ has finite length, the simple complex
$\int_2 I^{\bullet\bullet}$ is a complex of injective type, and the
morphism of complexes}

\VerdierSemanticAnchor{ch4d-4-6-8-morphism}
\[
\textit{$\int_2\varepsilon\colon
Y\longrightarrow\int_2 I^{\bullet\bullet}$}
\]
\textit{is a quasi-isomorphism.}

\VerdierSemanticAnchor{ch4d-4-6-8-proof-spectral-objects}
Indeed, the morphism $\varepsilon$ induces a morphism between the first
spectral object of $Y$ with values in $D(\mathcal A)$ (4.1.5) and the
spectral object with values in $D(\mathcal A)$ defined by
$I^{\bullet\bullet}$ filtered by its first degree ((4.1.7), (4.1.8)).
These two spectral objects are carried by the canonical cohomological
functor
$H\colon D(\mathcal A)\longrightarrow\mathcal A$ to spectral objects with
values in the abelian category $\mathcal A$. It suffices to show that
$\varepsilon$ induces an isomorphism between these spectral objects. Now
the spectral object defined by $Y$ is always stationary (Chapter~II,
4.4.2), and under the hypotheses made the spectral object defined by
$I^{\bullet\bullet}$ is stationary. It therefore suffices to show that
$\varepsilon$ induces an isomorphism on the initial terms. These initial
terms are, on the one hand,

\VerdierSemanticAnchor{ch4d-4-6-8-initial-y}
\[
\mathrm I_1^{p,q}=H^q(Y^p),
\]
and, on the other hand,

\VerdierSemanticAnchor{ch4d-4-6-8-initial-i}
\[
{}'\mathrm I_1^{p,q}=H^q(I^{p,\bullet}).
\]
Moreover, the morphism induced by $\varepsilon$ on the initial terms is
obtained by passing to cohomology in the morphism of complexes

\VerdierSemanticAnchor{ch4d-4-6-8-epsilon-p}
\[
\varepsilon^p\colon Y^p\longrightarrow I^{p,\bullet}.
\]

\VerdierSemanticAnchor{ch4d-4-6-8-proof-end}
But this morphism is a quasi-isomorphism, which proves the proposition.

\VerdierSemanticAnchor{ch4d-4-6-9}
\medskip\noindent\textbf{4.6.9.}\quad
Let $F\colon D(\mathcal A)\longrightarrow\mathcal B$ be a cohomological
functor, let $Y$ be an object of $\operatorname{comp}(\mathcal A)$, and
let $Y\longrightarrow I^{\bullet\bullet}$ be a Cartan--Eilenberg
injective resolution. The functor $F$ carries the corresponding
Cartan--Eilenberg spectral object (4.6.7) to a spectral object with values
in $\mathcal B$, and hence to a spectral sequence called the
\emph{Cartan--Eilenberg spectral sequence}. Denote by
$(\mathrm{C.E.})_r^{p,q}$ the terms of this spectral sequence
(Chapter~II, 4.3.3.1). The initial terms of this spectral sequence are
(4.1.9.1)

\VerdierSemanticAnchor{ch4d-4-6-9-initial-terms}
\[
(\mathrm{C.E.})_1^{p,q}
=F^{p+q}\bigl(I^{\bullet,p}[-p]\bigr)
=F^q(I^{\bullet,p}),
\]
\endgroup
